\documentclass[10pt,a4paper]{amsart}
\usepackage[margin=19mm]{geometry}
\usepackage{amsmath,amssymb,mathtools}
\usepackage[scr=rsfs,cal=euler]{mathalfa}
\usepackage{xfrac}
\usepackage[unicode,hidelinks,bookmarks=false]{hyperref}
\newcommand{\ie}{i.e.\ }
\newcommand{\cf}{cf.\ }
\newcommand{\resp}{respectively\ }
\newcommand{\Subsection}{\subsection}
\providecommand{\nobreakdash}{\nobreak\hbox{-}}
\setlength{\emergencystretch}{2em}
\multlinegap0pt
\usepackage{bm}
\usepackage{bbm}
\usepackage{dsfont}
\usepackage{xspace}
\usepackage{cancel}
\usepackage{mathtools}
\usepackage{amsthm}
\makeatletter
\@ifundefined{defn}{\newtheorem{defn}{Defn}}{}
\@ifundefined{lem}{\newtheorem{lem}[defn]{Lemma}}{}
\makeatother
\newcommand\PP{\mathbb P}
\newcommand\Pic{\operatorname{Pic}}
\newcommand\Z{\mathbb Z}
\newcommand\ord{\operatorname{ord}}
\title[A proof of N.Takahashi's conjecture for $(\mathbb{P}^2,E)$ a]{A proof of N.Takahashi's conjecture for $(\mathbb{P}^2,E)$ and a refined sheaves/Gromov-Witten correspondence}
\author{Pierrick Bousseau}
\date{Source: arXiv:1909.02992v3 (CC BY 4.0)}
\begin{document}
\maketitle

\noindent\emph{Source and licence.} A proof of N.Takahashi's conjecture for $(\mathbb{P}^2,E)$ and a refined sheaves/Gromov-Witten correspondence. Version arXiv:1909.02992v3 (CC BY 4.0), distributed under the Creative Commons Attribution 4.0 International licence, as recorded in the arXiv licence metadata for this exact version and shown on the abstract page https://arxiv.org/abs/1909.02992v3. Full rights notice: licenses/arxiv-1909.02992-CC-BY-4.0.txt.\par\smallskip

\noindent\textit{Editorial scope.} Counted unit: the Lem whose source label is \texttt{lem\_monodromy} in the article (source0.tex lines 1498--1502, source label \texttt{lem\_monodromy}), delivered together with the article's own complete original proof of it (source0.tex lines 1504--1528, one \texttt{proof} environment of 25 lines). No mathematics is written, repaired or re-derived: statement and proof are the article's own text line for line. Required context retained uncounted: lem\_divisibility (lines 1475--1482). Every definition, display and label of those lines is retained unchanged. Omitted from this package and not redistributed: 1--1474, 1483--1497, 1503--1503, 1529--3699. Nothing inside a retained excerpt is omitted or altered: every retained body is delivered exactly as the article's lines are written, so no omission receipt of any kind accompanies this package. Citations: the retained excerpts carry no citation command at all, so no bibliography entry of the article is transcribed and no reference is redistributed. Reference adaptation: none was needed and none was made; every reference inside the delivered unit resolves inside it, so no unresolved reference or citation is produced. Printed slips kept literal: no printed word, symbol, hypothesis or number of the counted statement or of its proof was altered. Layout and class: the article's own class is \texttt{amsart} with packages including babel, inputenc, amsmath, graphicx, todonotes, amssymb, amscd, latexsym, epsfig, xy, euscript, amsthm, tikz, MnSymbol; the wrapper is the lane's pinned \texttt{amsart} editorial wrapper at 10pt on A4, which restates the theorem-like environments the retained text needs and adds the article's own macro definitions that the retained text needs (\texttt{\textbackslash PP}, \texttt{\textbackslash Pic}, \texttt{\textbackslash Z}, \texttt{\textbackslash ord}), with extra wrapper packages (bm, bbm, dsfont, xspace, cancel, mathtools). The delivered numbering is the unit's own. Assets: no retained span contains a figure environment, a graphics-inclusion command, a PDF-page inclusion, a table or a TikZ picture; no image file is copied into this package, the package declares no asset dependency and no image file is redistributed with it. Licence: version arXiv:1909.02992v3 of this article is distributed under the Creative Commons Attribution 4.0 International licence, as recorded in the arXiv licence metadata for this exact version and shown on the abstract page https://arxiv.org/abs/1909.02992v3. A full rights notice is delivered at licenses/arxiv-1909.02992-CC-BY-4.0.txt. Characters: every retained excerpt is pure ASCII, so the delivered engine, XeLaTeX with its default Latin Modern text font, reports no missing character.
\begin{lem}\label{lem_divisibility}
Let $d \in \Z_{>0}$, $p \in P_d$, and let $p_0$ be a flex point of $E$.
\begin{enumerate}
    \item If $3$ divides $d(p)$, then $\ord(p-p_0)=3d(p)$.
    \item If $3$ does not divide $d(p)$, then 
    $\ord(p-p_0)=3d(p)$ or $\ord(p-p_0)=d(p)$.
\end{enumerate}
\end{lem}

\begin{lem} \label{lem_monodromy}
For every $d \in \Z_{>0}$
and $g\in \Z_{\geqslant 0}$, the 
Gromov-Witten invariants $N_{g,d}^{\PP^2/E,p}$ only depend on the point $p \in P_d$ through the integer $d(p)$. 
\end{lem}

\begin{proof}
We fix a $d \in \Z_{>0}$ and two points $p, p' \in P_d$
such that $d(p)=d(p')$. We have to show that 
$N_{g,d}^{\PP^2/E,p}=N_{g,d}^{\PP^2/E,p'}$.

We show below that there exists a flex point $p_0$ of $E$ such that $\ord(p-p_0)=\ord(p'-p_0)$.
This will be enough. Indeed, the monodromy of the family of smooth cubics $E$ in $\PP^2$ maps surjectively to $SL(2,\Z)$
acting on $\Pic^0(E)[3d] \simeq (\Z/(3d))^2$, and so if 
$\ord(p-p_0)=\ord(p'-p_0)$, then $p$ and $p'$ can be related by monodromy and so the result follows from deformation invariance in relative Gromov-Witten theory.

We first choose $p_0$ any flex point of $E$. If $3$ divides 
$d(p)=d(p')$, then by Lemma 
\ref{lem_divisibility}, $\ord(p-p_0)=3d(p)=3d(p')=\ord(p-p_0)$ and we are done.

If $3$ does not divide $d(p)=d(p')$ and $\ord(p-p_0)=\ord(p'-p_0)$, then we are also done.
Therefore, by Lemma \ref{lem_divisibility}, and up to exchanging $p$ and $p'$, we can assume that $3$ does not divide $d(p)=d(p')$,
$\ord(p-p_0)=3d(p)$ and $\ord(p'-p_0)=d(p)$.
Then $d(p)(p-p_0)$ is nonzero $3$-torsion.
If $t$ is nonzero $3$-torsion, then 
$d(p)t$ is also nonzero $3$-torsion as $3$ does not divide 
$d(p)$. So if $t_1$ and $t_2$ are both nonzero $3$-torsion,
with $d(p)t_1=d(p)t_2$, then $t_1=t_2$. It follows that there exists $t$ nonzero $3$-torsion such that $d(p)t \neq -d(p)(p-p_0)$, and so $\ord(p-p_0+t)=3d(p)$ and 
$\ord(p'-p_0+t)=3d(p)$. So it is enough to replace 
$p_0$ by $p_0-t$.
\end{proof}

\end{document}
